\begin{claim}{B.40}{Proved}{Complete annulus and M\"obius-band quotients in genus three}
There are exactly $57$ mixed finite genus-three action classes whose
full quotient is an annulus or a M\"obius band with its cone and mirror
corner structure.  There are $39$ annulus classes and $18$ M\"obius-band
classes.  Their positive quotient is a torus.  Every group, generating
vector, quotient signature and complete section poset is given in
\path{type_3_0_torus_mirrors_catalogue.csv},
\path{type_3_0_torus_mirrors_certificates.json}, and
\path{type_3_0_torus_mirrors_sections.json}.
\end{claim}
\begin{proof}
An involution reversing the orientation of a torus has annulus,
M\"obius-band or Klein-bottle quotient.  Indeed its fixed set consists
of disjoint circles, and its quotient has Euler characteristic zero;
the orientation double is the torus.  The classification of compact
surfaces leaves the orientable annulus, the nonorientable M\"obius
band, and the closed nonorientable Klein bottle.  Their orientation
doubles give respectively two, one, and zero fixed circles.  Cone
points on a fixed circle become mirror corners, and other cone points
are paired to interior cones.  The Klein-bottle cases are B.41.

The complete positive input of B.33 leaves the following seven
signature/group pairs.  In each row $r$ is the number of positive
quotient cone points, all of order $m$; the last two columns are the
executed numbers of full annulus and M\"obius-band classes.
\[
\begin{array}{c|c|r|r|r|r}
 \text{positive row}&H&m&r&\text{annulus}&\text{M\"obius}\\\hline
 \mathrm O3.3&C_2&2&4&10&3\\
 \mathrm O3.5&C_3&3&2&5&4\\
 \mathrm O3.7&C_4&2&2&8&4\\
 \mathrm O3.8&C_2^2&2&2&9&4\\
 \mathrm O3.12&S_3&3&1&3&1\\
 \mathrm O3.17&D_8&2&1&3&1\\
 \mathrm O3.18&Q_8&2&1&1&1
\end{array}                                             \tag{B.40.1}
\]
Here $D_8$ has order eight.  For $l$ interior cones, the total number
of corners is $r-2l$.  On an annulus distribute them as
$k_1\le k_2$ with $k_1+k_2=r-2l$; on a M\"obius band place all
$k=r-2l$ on its single boundary.  All their orders are $m$, so these
integer choices already classify the colored quotient signatures.

We first give an exhaustive construction of the possible oriented
finite groups, rather than assuming a small-group name list.
Choose a negative coset representative $t$ in an index-two extension
of $H$.  Let $\alpha(h)=t h t^{-1}$ and $A=t^2$.  Its complete
extension data satisfy
\[
              \alpha^2=\operatorname{Inn}(A),\qquad
                          \alpha(A)=A.                 \tag{B.40.2}
\]
Conversely the product of (B.39.4), with these data, defines that
extension:
\[
 (h,\epsilon)(k,\delta)=
       (h\alpha^\epsilon(k)A^{\epsilon\delta},
                                  \epsilon+\delta\bmod2).
                                                        \tag{B.40.3}
\]
Every oriented group isomorphism restricts to some $\beta\in
\operatorname{Aut}(H)$ and sends $t$ to some $k t'$ in the negative
coset.  Therefore the entire isomorphism action on (B.40.2) is generated
by transport by $\beta$ and by
\[
       (\alpha,A)\longmapsto
       \bigl(\operatorname{Inn}(k)\alpha,\ k\alpha(k)A\bigr).
                                                        \tag{B.40.4}
\]
The converse follows by sending the displayed generators accordingly
in (B.40.3), so this is an exact fusion test.  The program enumerates
$\operatorname{Aut}(H)$ by assigning images to a generating set and
propagating on its complete Cayley graph, enumerates every pair
(B.40.2), and takes precisely the finite orbits (B.40.4).
For the seven $H$'s in (B.40.1), it obtains respectively
$2,2,4,3,1,4,4$ oriented extension classes.

Here are the full source presentations and enumeration conditions.
For an annulus, put $X=x_1\cdots x_l$ and use positive generators
$x_i,e_1,e_2$ with
\[
                x_i^m=1,\qquad X e_1e_2=1.             \tag{B.40.5}
\]
On its $j$th boundary there are reflections
$c_{j,0},\ldots,c_{j,k_j}$ satisfying
\[
 c_{j,s}^2=1,\quad(c_{j,s-1}c_{j,s})^m=1,\qquad
               c_{j,k_j}=e_j c_{j,0}e_j^{-1}.          \tag{B.40.6}
\]
For $k_j=0$, there is one reflection commuting with $e_j$.
For a M\"obius band there is one negative crosscap generator $d$,
one boundary with (B.40.6), and the relation
\[
                X e d^2=1,
                   \qquad e=X^{-1}d^{-2}.              \tag{B.40.7}
\]
In each finite extension the program assigns every possible positive
order-$m$ element to every $x_i$, every positive element to $e_1$
in the annulus, and every negative element to $d$ in the M\"obius
band.  It then assigns every sequence of negative involutions to
the mirrors, requiring exact order $m$ for each adjacent product
and the stated closing relation.  It retains precisely the vectors
generating the whole finite group.  Two involutions whose product
has exact order $m$ generate a dihedral group of order $2m$; thus
all local groups inject.  The orientation character takes the
specified signs on every generator, so the kernel is oriented and
contains no reflection.  The signature equation gives
$|K|\chi^{\rm orb}(Q)=-4$; its torsion-free surface kernel therefore
has genus three.  Conversely every action with this quotient yields
one of these enumerated epimorphisms.

Target automorphisms are removed by the exact rooted colored Cayley
graph test of B.37, now applied to the full NEC generating vector.
Its prescribed generator signs force any resulting isomorphism to
preserve the orientation character.  For source fusion, the following
are explicit complete generators.  The interior braids are
(B.37.7), with all other generators fixed.  A cyclic boundary move
sends $c_{j,s}$ to $c_{j,s+1}$, extending the reflection sequence
by (B.40.6), while fixing $x_i$, $e_j$, and the crosscap generator.
After $k_j$ moves it conjugates that boundary's reflection sequence
by $e_j$ and fixes the other generators, which supplies the boundary
Dehn twist where it is needed.

For the annulus, pushing the last interior point once around the
first boundary, with $a=x_l$ and $e=e_1$, has
\[
 \begin{aligned}
 a&\longmapsto (a e)a(a e)^{-1},\\
 e_1&\longmapsto a e_1a^{-1},\\
 c_{1,s}&\longmapsto a c_{1,s}a^{-1}.
 \end{aligned}                                           \tag{B.40.8}
\]
It fixes the other $x_i$, $e_2$ and the second reflection sequence.
This is the square of the Hurwitz exchange between the puncture
meridian and the first boundary loop; in particular it preserves
$a e_1$ exactly.  The interior braids conjugate this to pushes of
each other interior point.  When $k_1=k_2$, interchange of the two
boundary components is
\[
 \begin{aligned}
 e_1&\longmapsto e_1e_2e_1^{-1},&
 e_2&\longmapsto e_1,\\
 c_{1,s}&\longmapsto e_1c_{2,s}e_1^{-1},&
 c_{2,s}&\longmapsto c_{1,s},
 \end{aligned}                                           \tag{B.40.9}
\]
fixing the interior generators.  This is their orientation-preserving
boundary exchange, with the chosen connecting paths.

The remaining annulus generator reverses the orientation of its
fundamental region.  For the peripheral list
$(w_1,\ldots,w_{l+2})=(x_1,\ldots,x_l,e_1,e_2)$, write
$W_i=w_1\cdots w_i$.  It sends
\[
             w_i\longmapsto W_{i-1}w_i^{-1}W_{i-1}^{-1}.
                                                        \tag{B.40.10}
\]
For the first boundary its reflection sequence is reversed and
conjugated by $X$; for the second it is reversed and conjugated by
$X e_1$.  Reversing a cycle means
$(c_0,\ldots,c_{k-1})\mapsto(c_k,c_{k-1},\ldots,c_1)$,
with $c_k=e c_0e^{-1}$; an empty cycle retains its single reflection.
These substitutions verify (B.40.5)--(B.40.6) directly.

For the M\"obius band, a one-sided slide of the last interior point
is particularly explicit.  Set
\[
          a=x_l,\qquad a'=(a d)a^{-1}(a d)^{-1}.
\]
The substitution is
\[
 \begin{aligned}
  x_l&\longmapsto a',& d&\longmapsto a d,\\
  c_s&\longmapsto a'^{-1}a\,c_s\,a^{-1}a'.
 \end{aligned}                                           \tag{B.40.11}
\]
It fixes the other interior generators and determines $e$ by
(B.40.7).  To obtain it geometrically, regard the boundary temporarily
as the last peripheral loop in the projective-plane presentation:
exchange $(a,e)$ by its Hurwitz move, slide the now-last puncture
by (B.39.5), and undo that exchange.  This gives exactly
(B.40.11), including the displayed transport of every boundary
reflection.  Finally a reversal of the core has
\[
 \begin{aligned}
 (x_1,\ldots,x_l)&\longmapsto(x_l^{-1},\ldots,x_1^{-1}),\\
 d&\longmapsto d^{-1},\qquad
 c_s\longmapsto X c_{k-s}X^{-1}.
 \end{aligned}                                           \tag{B.40.12}
\]
Here $e$ goes to $X e^{-1}X^{-1}$, so the closing relation is
preserved.

We justify that this is the whole source action.  First prescribe
the cyclic permutations, and possible reversal or interchange, of
the marked boundary components.  The boundary moves above realize
exactly these permutations.  Upon forgetting the interior points,
a remaining annulus map is generated by its core twist, and a
remaining M\"obius-band map by reversal of its core.  Cutting along
a spanning arc gives this classification: the annulus permits the
integer winding of the arc, while cutting the M\"obius band reduces
a boundary-fixed map to a disc isotopy.  The boundary reversal is
(B.40.12).  The core twist is already among the full cyclic boundary
moves when both boundaries are marked; when a boundary is unmarked
its rotation isotopes it away in the unpunctured quotient.
The kernel of forgetting interior points consists of their braids.
An interior point in an annulus minus the other points has fundamental
group generated by the puncture meridians and the core loop; in a
M\"obius band the corresponding extra generator is one-sided.
Forgetting points one at a time, and using isotopy extension, therefore
shows that their braid groups are generated by the ordinary interior
braids and respectively (B.40.8) or (B.40.11).  This also includes
any twist that becomes trivial only after the interior points are
forgotten.  Thus every geometric source automorphism is included,
and each included automorphism is realized by the maps just described.
The kernel-to-cover lifting argument of B.37 proves that the computed
orbits are exactly the surface conjugacy classes.

Execution gives the $57$ classes in (B.40.1).  There is an independent
check on the source generators in every one-cone, no-corner case:
the complete primary automorphism lists in
\cite[Lemmas 3.7 and 3.8]{C_GromadzkiSzepietowski2016} give exactly the
same partitions of the entire target-normalized epimorphism sets.
The file \path{check_torus_mirror_generators_output.json}
records all six comparisons, comparing full orbit memberships and
not just their numbers.

The program then computes all section invariants.  Its mirror
incidence graph has a separate closed cycle for each quotient
boundary, with endpoint gluing by that boundary's actual element
$e_j$.  At every corner it uses the full local dihedral group, and
checks that the fixed subgraph of each reflection has degree two
at every vertex.  Consequently its connected components are exactly
the fixed ovals, including all their intersections with other
reflection fixed sets.  Cone and corner cosets give the isolated
fixed points.  The procedure of B.26 then supplies every section,
multiplier, centralizer, extension incidence, maximal section, group
orbit and stabilizer.  All $57$ presentations, local injections,
Euler characteristics and complete section calculations passed.
\end{proof}

The $57$ classes have orders $4,6,8,12,16$ with multiplicities
$13,9,25,4,6$, respectively.  The full quotient surfaces in this
claim have boundary, whereas those in B.41 are closed.  A conjugacy
of actions descends to a homeomorphism of their full quotient
orbifolds and therefore preserves boundary.  Their positive
quotients have genus one, whereas those in B.37 and B.39 have genus
zero and those in B.38 have genus two or three.  The induced
homeomorphism of positive quotients preserves genus.  These exact
quotient maps prove that the four portions contribute disjoint
conjugacy classes.
